\documentclass[11pt]{article}
\usepackage{amsmath,hyperref,geometry}
\geometry{margin=1in}
\title{Warehouse Lane Cutover When an Independent Owner Withdraws}
\date{}
\begin{document}
\maketitle

\paragraph{Claim.}
A lawful owner can stop proposing work before the shared actuator stops receiving its prior proposals. A safe cutover therefore requires exclusion of old starts and actual drainage of old physical use, or an independent path serving the new duty without declaring the old one discharged. A signed statement from the withdrawing owner does not provide either fact by itself.

\paragraph{Authority, promises, and phases.}
Two independent maintenance crews share a unit-capacity indoor warehouse lane. A site-owned trusted port mediates every physical start, including starts by a crew that does not follow the agreement. The crews cannot bypass the port or change its sensors, calendar, or phase order. An old owner may refuse new proposals. It can withdraw its own future participation at tick 0, when the site sends the cutover instruction. A proposal sent before tick 0 can reach the port at tick $a$. It is not a firm advance booking: it becomes a noncancelable service commitment exactly on port acceptance. The agreement honors a valid pending proposal if it arrives while the old-start port is open and the lane vacant. The port fence takes effect at tick $f$, $1\le f\le L$, and prohibits all later old starts, but neither cancels prior commitments nor accelerates their completion. The firm old job occupies the lane on $[a,a+d)$ with $1\le d\le D$ in the bounded service envelope. The successor owes the customer a start no later than $S$ (completion one tick later), and cannot refuse admitted work to obtain a passing score. Physical release happens before the fence, then applicable old arrivals and receipt delivery, then successor admission at the same tick. Fence time is visible at the shared port without extra acknowledgment latency. One old proposal suffices for all sharp lower witnesses; the formula still bounds a larger accepted set if the same serial resource and per-job dwell bounds hold. Only the one-proposal case is replayed.

\paragraph{Legal withdrawal, accident, arbitrary deviation.}
Refusal to propose optional work creates no inherited duty and is not a fault. In the lawful matched pair the same proposal sent at $-2$ is delivered at 1 versus 2, with $f=2$. It starts in the first case and is rejected in the second, with identical withdrawal notices. In the first, $d=3$ means the logical revocation at 2 leaves physical use in place until 4. The identity of the accepted job, not an owner-wide handoff label, distinguishes the obligations. If a proposal constituted a binding upstream booking, rejecting it at the fence would not suffice: the site must lawfully rebook it with equivalent service, honor it via independently usable capacity, or obtain authorized release. None of these may be inferred from the fence. This stronger booking model is not counted as served here.

An accidental jam can prevent a previously started job from physically releasing. Safety remains conditional on the site interlock, but the jam is outside the bounded old-service envelope and physical completion is no longer promised. The replay has one non-releasing jam, with and without an independently held compatible support. The reserve can meet the successor promise but cannot settle the old duty; that remains outstanding. An arbitrary deviation is modeled as an unreliable owner-side vacancy declaration. The old owner's identity does not make this declaration evidence. The documentary policy uses only a site-authenticated physical release record, and the local policy uses the physical interlock, not the declaration. Even noncompliant requests cannot change admission beyond the fence in this model; no guarantee is made against bypass of the mediating port, corruption of the independent sensor, or withholding a necessary physical release. The binding rules, not the labels assigned to the owners, govern the claims.

\paragraph{Exact dispatch and necessity.}
With direct sensing the successor can start on the shared lane exactly on the first tick at which the port has fenced old starts and the lane is vacant. For an accepted old job this is $\max(f,a+d)$; otherwise it is $f$. Vacancy without exclusion is insufficient: before $f$ an in-flight proposal may still arrive, and either accepting it risks a conflict or refusing it changes the incumbent promise. This conclusion is for the original interface, not a different actuator that atomically fences and acquires at an earlier tick. Exclusion without vacancy is insufficient by physical exclusivity; an occupied lane cannot be declared free without changing the plant. A vacancy sample obtained before the fence cannot distinguish a later pre-fence admission from no admission. At tick 2 in the matched pair, the same owned declaration can accompany opposite true occupancy states. A proof based solely on that declaration must either decline the prompt start in the empty case or authorize a potentially conflicting start in the occupied case. The trusted interlock still blocks a physically impossible order; blocking is not delivery of the service promise.

The worst shared-lane start is
\[
 W_{\rm physical}=\max(L,L-1+D).
\]
The latest lawful old start is $L-1$, attained at fence $L$; empty histories can still await a fence as late as $L$. Every accepted duty completes at its own $a+d$; rejected revocable proposals are not completions, and held duties cannot be erased. Thus $S\ge W_{\rm physical}$ is necessary and sufficient for a universal shared-lane start under the declared law. A certified shared clock/dwell policy after the actual fence can use $\max(f,f-1+D)$ even without a vacancy sensor. It attains the same worst bound but need not give the earliest safe start on each history. For $S<W_{\rm physical}$ the witness rules out relying on that lane under the same non-preemption and authority law; a separately enforceable service path or a legitimate promise change is required. Unbounded occupation admits no finite shared-lane service guarantee. An independent spare lane can serve the successor without waiting, but consumes a separate held support and neither shortens old occupation nor verifies settlement.

\paragraph{What documents can prove.}
The documentary gate knows the complete accepted set at fence time (including an affirmative empty-set attestation) and waits for a trusted, identity-bound receipt for each member. Acknowledging fence or observing an issuer's list is not evidence that its list covers physically accepted work. Receipt is produced at actual release and arrives at most $R$ ticks later. The rule's tight worst time is $W_{\rm receipt}=\max(L,L-1+D+R)$. At $(L,D,R)=(2,3,2)$ the old job can still be present at 3, clears at 4 and is documented at 6; only the direct or certified-timeout gate meets $S=4$ on that history. At tick 3, a two-tick old job and a three-tick old job can expose identical documentary views (same fenced identity, neither receipt available) but opposite physical vacancy. This is a genuine local-information limit for prompt release with those observations. The additive $R$ is a cost of the particular affirmative-receipt rule, not a lower bound against a known dwell-and-clock bound. If evidence generation itself is delayed, the formula must include that delay; if the complete set or its receipts can be withheld indefinitely, documentary liveness is no longer established.

\paragraph{Comparators and evidence.}
\href{https://static.usenix.org/events/osdi06/tech/full_papers/burrows/burrows_html/}{Burrows, Chubby, \S2.4} already places sequencer checking at the recipient, not just the owner. Here the recipient is a nonbypassable physical port; sequencer-style checking does not dispose of an occupied lane. \href{https://www.cs.utexas.edu/~aim/papers/JAIR08-dresner.pdf}{Dresner and Stone, \S3.2} have distinct reservation cancellation and identity-bearing completion messages; their Done receipt is not the trusted mandatory physical settlement signal stipulated here. A conventional direct proof combines recipient fencing with a drain guard (physical, bounded time, or affirmative receipts), obtaining exactly the same service and safety bounds. No Byzantine Control theorem is imported: fault taxonomy motivates the distinctions, but there is no quorum or owner-trust conclusion under changed authority, phase, physical, or service assumptions. The premises and guarantees are for this modeled civil resource, not an operational claim.

The event replay uses phased starts, releases and arrivals. A separate interval oracle scans permissible dispatch times and checks 22,500 policy results on 7,500 parameterized traces, including starts before withdrawal. The exact bounds follow from the argument above, independent of the sweep. The featured traces include lawful refusal, the two reached withdrawal cuts, documentary delay, a false owner declaration blocked by the interlock, and a physically stuck duty with and without independent capacity. These are finite checks and written proofs, not physical or formal verification.
\end{document}
